\uuid{tOjI}
\titre{Matrice d'une application linéaire}
\auteur{hanani ; mzari}
\organisation{exo7}
\licence{CC BY-NC-SA 4.0}
\datecreate{2020-01-20}
\niveau{L1}
\module{Algèbre}
\chapitre{Matrice}
\sousChapitre{Matrice et application linéaire}
\difficulte{1}
\theme{Matrice d'une application linéaire}
\exo7qcmid{lille-2.matrice-application.003}

\contenu{
\begin{question}
\qclassification{108.03}
\qtags{motcle=Matrice d'une application linéaire}
On considère $\Rr^2$ et $\Rr^3$ munis de leurs bases canoniques et $f$ l'application linéaire définie par :
$$\begin{array}{rccc}f:&\Rr^2 &\to& \Rr^3\\
& (x,y)&\to &(y,x,-y).  \end{array}$$
La matrice de $f$ relativement aux bases canoniques est :
\begin{answers}
\bad{$\left(\begin{array}{rcc}
0&1&0\\
1&0&-1\end{array}\right)$.}
\bad{$\left(\begin{array}{rcc}
0&1&0\\
1&0&-1\\
0&0&0\end{array}\right)$.}
\good{$\left(\begin{array}{rc}
0&1\\1&0\\0&-1\end{array}\right)$.}
\bad{$\left(\begin{array}{rcc}
0&1&0\\
1&0&0\\0&-1&0\end{array}\right)$.}
\end{answers}
\begin{explanations} Soit ${\cal {B}}=\{e_1,e_2\}$  et ${\cal {B}}'=\{e_1',e_2',e_3'\}$ les bases canoniques de  $\Rr^2$ et  $\Rr^3$ respectivement. La matrice de $f$  relativement à ces bases est la matrice dont la $1$ère colonne $f(e_1)=\left(\begin{array}{r}
0\\1\\0\\\end{array}\right)$ et la 2ème colonne est $f(e_2)=\left(\begin{array}{r}1\\0\\-1\\
\end{array}\right)$. Cette matrice est :
 $\left(\begin{array}{rc}0&1\\
1&0\\0&-1\end{array}\right).$
\end{explanations}
\end{question}
}
